\documentclass[12pt,a4paper]{article}

\usepackage[utf8]{inputenc}
\usepackage[T1]{fontenc}
\usepackage{lmodern}
\usepackage[margin=1in]{geometry}
\usepackage{booktabs}
\usepackage{tabularx}
\usepackage{longtable}
\usepackage{array}
\usepackage{hyperref}
\usepackage{amsmath}
\usepackage{graphicx}
\usepackage{xcolor}
\usepackage{caption}
\usepackage{setspace}

\hypersetup{
    colorlinks=true,
    linkcolor=blue,
    citecolor=blue,
    urlcolor=blue
}

\newcolumntype{C}[1]{>{\centering\arraybackslash}p{#1}}

\begin{document}
\title{\textbf{Indonesia's Tariff Agreement After the US Supreme Court Ruling: What It Means for Indonesia}}


\author{Teguh Yudo Wicaksono\thanks{Faculty of Economics and Business, UIII} \and
        Rima P Artha \and
        Philips J.\ Vermonte\thanks{Faculty of Social Science, UIII} \\[0.5em]
        \small Indonesian Institute for Foreign Affairs (IIFA),\\
        \small Universitas Islam Internasional Indonesia (UIII)}


\date{February 28, 2026}
\onehalfspacing
\maketitle

%----------------------------------------------------------------------
\section*{Summary}
%----------------------------------------------------------------------

On February 19, 2026, Indonesia and the United States signed the Bilateral Agreement of Reciprocal Trade (ART), locking in a 19\% reciprocal tariff and granting zero-tariff exemptions to 1,819 product lines---including palm oil, rubber, coffee, and cocoa. The ink was barely dry. The next day, the US Supreme Court struck down the entire legal basis of the agreement, ruling 6--3 in \textit{Learning Resources, Inc.\ v.\ Trump} that the International Emergency Economic Powers Act (IEEPA) does not authorize the president to impose tariffs.

Within hours, the Trump administration invoked Section 122 of the Trade Act of 1974---a rarely used balance-of-payments provision---at 10\%, then escalated to \textbf{15\%}, the statute's maximum, on February 21. Unlike IEEPA, Section 122 cannot differentiate by country and carries a mandatory 150-day expiry. The bilateral deal Indonesia had spent months negotiating no longer has a legal home.

Contrary to initial reports, Indonesia's \textbf{average tariff rate under Section 122 at 15\% (19.0\%) is slightly \textit{below} the deal rate (19.7\%)}---not above it. The reason is that the Section 122 proclamation carries Annex II product exemptions that mirror the IEEPA Annex II list, covering coffee, cocoa, and partial exemptions for electronics and rubber. Once these exemptions are applied alongside updated 2025 trade data (\$34.84 billion, up 24\% on pre-tariff front-loading), the corrected aggregate average tariff rate is 19.0\% under S122 at 15\%---0.7 percentage points below the deal.

The distribution of burden, however, is sharply unequal. Palm oil and rubber exporters are the clear losers: stripped of their deal exemptions and not covered by the Annex II, they face tariff increases of 11--15 percentage points relative to what the deal would have provided. Coffee and cocoa exporters, by contrast, are largely unaffected---both sectors are fully Annex II exempt under S122, meaning they pay only their low Most-Favored-Nation (MFN) rates as they would have under the deal. Manufacturing exporters---apparel, footwear, and furniture---see modest tariff reductions of 2--4 percentage points compared to the deal's rate. Indonesia's concessions under the ART, including \$33 billion in purchase commitments and a 20-year extension of PT Freeport Indonesia's mining license, were made in exchange for tariff terms that can no longer be legally delivered.

\newpage
%----------------------------------------------------------------------
\section{Introduction}
%----------------------------------------------------------------------

The recent development on the tariffs imposed by President Trump on the US' trading partners has raised controversies---not only in the targeted countries, but also within the United States itself. Domestic political battles between President Trump and the Supreme Court of the United States (SCOTUS) in the judicial branch continue to reshape how trade deals signed with many countries, Indonesia included, will actually take effect.

On April 2, 2025, invoking his interpretation of the International Emergency Economic Powers Act (IEEPA), President Trump announced the so-called ``Liberation Day'' Executive Order that imposed a 32\% reciprocal tariff on Indonesia. A 90-day pause was offered to countries willing to negotiate. By late July 2025, Indonesian and US negotiators reached a framework agreement reducing the rate to 19\%.

Presidents Prabowo Subianto and Donald Trump signed the Bilateral Agreement of Reciprocal Trade (ART) on February 19, 2026.\footnote{President Trump also utilizes Section 232 of the Trade Expansion Act of 1962 and Section 301 of the Trade Act of 1974.} The following day, SCOTUS invalidated the IEEPA tariff authority. In response, President Trump implemented a Section 122 tariff of 10\%---raised to 15\% within hours.

The SCOTUS decision can be read through two lenses. First, it reflects the constitutional system of checks and balances functioning as intended: the judiciary restraining executive overreach on a power---the authority to impose taxes and duties---that the Constitution grants to Congress alone. Second, for Indonesia, the practical consequences of this democratic balancing act are economically significant. A months-long negotiation produced a deal with product-specific carve-outs carefully matched to Indonesia's export profile; a flat emergency tariff applied uniformly to all countries produces a very different distribution of burden.

Countries that signed bilateral agreements under now-invalidated IEEPA authority face a fundamental question: what do those concessions mean when the US cannot legally deliver the tariff terms they bought? For Indonesia, this analysis calculates the concrete effects on the tariff burden facing its exporters, identifies which sectors gain and lose relative to the deal, and outlines the policy options ahead.

%----------------------------------------------------------------------
\section{Context: The US--Indonesia Trade Relationship}
%----------------------------------------------------------------------

Indonesia is America's 16th-largest goods trading partner. In 2025, the US imported \$34.84 billion from Indonesia---a 24\% increase from 2024's \$28.05 billion, largely reflecting pre-tariff front-loading by Indonesian and US importers anticipating higher duties \cite{census2024,ustr2024}. US exports to Indonesia in 2024 were approximately \$10.16 billion, leaving a goods trade deficit of around \$17.9 billion.

Indonesia holds dominant US market positions in several products: refined palm oil (85.1\% of US imports), industrial oleic acid (80.6\%), cinnamon (80.1\%), technically specified natural rubber (48\%), and crabs (46.8\%) \cite{comtrade}. Table~\ref{tab:topexports} shows Indonesia's top export sectors based on 2025 US import data, which reflect the surge in electronics and other goods ahead of the tariff changes \cite{usitc_dataweb}.

\begin{table}[h!]
\centering
\caption{Top Indonesian Export Sectors to the US (2025)}
\begin{tabular}{clccc}
\toprule
\textbf{Rank} & \textbf{Sector} & \textbf{HS Chapters} & \textbf{Value (\$B)} & \textbf{Share} \\
\midrule
1 & Electrical/Electronics & 85 & 7.84 & 22.5\% \\
2 & Footwear & 64 & 2.97 & 8.5\% \\
3 & Palm Oil/Fats & 15 & 2.32 & 6.7\% \\
4 & Knitted Apparel & 61 & 2.52 & 7.2\% \\
5 & Woven Apparel & 62 & 2.21 & 6.3\% \\
6 & Machinery & 84 & 1.65 & 4.7\% \\
7 & Furniture & 94 & 1.55 & 4.5\% \\
8 & Cocoa & 18 & 0.77 & 2.2\% \\
9 & Rubber & 40 & 1.71 & 4.9\% \\
10 & Fish Preparations & 16 & 0.77 & 2.2\% \\
\bottomrule
\end{tabular}
\par\smallskip
\begin{minipage}{\linewidth}
\footnotesize
\textit{Note:} Values are customs value of U.S.\ general imports from Indonesia, calendar year 2025 (12-month sum), aggregated at the HTS 2-digit chapter level. Share is each sector's value as a proportion of total U.S.\ imports from Indonesia (\$34.84 billion). The 2025 surge relative to 2024 reflects pre-tariff front-loading.\\
\textit{Source:} U.S.\ Census Bureau monthly trade statistics, January--December 2025 \cite{census2024}.
\end{minipage}
\label{tab:topexports}
\end{table}

Table~\ref{tab:chronology} traces the rapid sequence of policy changes that have defined the US--Indonesia trade relationship since April 2025. The pace has been unprecedented: in less than a year, Indonesia's tariff environment has oscillated from 5\% (MFN baseline) to 36.6\% (Liberation Day), down to 19.7\% (the negotiated deal), and now sits at 19.0\% under Section 122 with Annex II exemptions---all within the space of ten months.

\begin{table}[htbp]
\centering
\caption{Chronology of the Tariff Crisis}
\begin{tabularx}{\textwidth}{lX}
\toprule
\textbf{Date} & \textbf{Event} \\
\midrule
April 2, 2025 & ``Liberation Day'' executive order: 32\% reciprocal tariff on Indonesia via IEEPA \\
April 9, 2025 & 90-day pause announced for countries entering negotiations \\
June 2025 & Section 232 tariffs on steel/aluminum doubled from 25\% to 50\% \cite{crs_232} \\
July 7, 2025 & Trump letter threatening reimposition of 32\% if no deal \\
July 22, 2025 & US--Indonesia framework agreement: reduced to 19\% \\
August 1, 2025 & 19\% reciprocal tariff takes effect \\
February 19, 2026 & Bilateral Agreement of Reciprocal Trade (ART) finalized \cite{art2026}; 1,819 tariff lines exempted; Freeport MoU signed \cite{freeport_mou} \\
February 20, 2026 & SCOTUS strikes down IEEPA tariffs 6--3 in \textit{Learning Resources v.\ Trump} \cite{scotus2026}; Trump invokes Section 122 at 10\% \\
February 21, 2026 & Trump raises Section 122 tariff to 15\%---the statutory maximum \cite{s122_statute} \\
\bottomrule
\end{tabularx}
\par\smallskip
\begin{minipage}{\linewidth}
\footnotesize
\textit{Note:} Dates reflect the effective date of each executive action or announcement. All tariff rates are ad valorem.\\
\textit{Source:} Compiled from Executive Orders, White House fact sheets, \cite{scotus2026}, \cite{art2026}, \cite{freeport_mou}, \cite{crs_232}, and \cite{s122_statute}.
\end{minipage}
\label{tab:chronology}
\end{table}

Following the SCOTUS decision, all IEEPA reciprocal tariffs---including country-specific bilateral deals---were repealed. President Trump maintained Section 232 tariffs on steel, aluminum, and automobiles, and invoked Section 122 as an emergency replacement. Table~\ref{tab:legal survive} summarizes which authorities survived the ruling and which fell.

\begin{table}[htbp]
\centering
\caption{Legal Authorities: What Survives vs.\ What Falls}
\begin{tabularx}{\textwidth}{XX}
\toprule
\textbf{Survives} & \textbf{Falls} \\
\midrule
Section 232 tariffs (steel, aluminum, autos) & All IEEPA reciprocal tariffs \\
Section 301 tariffs (China) & IEEPA universal baseline tariff \\
Anti-dumping / countervailing duties & Country-specific bilateral deals under IEEPA \\
Tariff Act of 1930 duties & \\
Section 122 (invoked at 10\%, raised to 15\%) & \\
\bottomrule
\end{tabularx}
\par\smallskip
\begin{minipage}{\linewidth}
\footnotesize
\textit{Note:} Classification reflects the immediate legal effect of the Supreme Court's 6--3 decision in \textit{Learning Resources, Inc.\ v.\ Trump} \cite{scotus2026} on February 20, 2026. Section 122 was invoked by executive order on the day of the ruling and raised to the 15\% statutory maximum on February 21, 2026.\\
\textit{Source:} Authors' analysis based on \cite{scotus2026}, \cite{s122_statute}, \cite{crs_232}, and \cite{crs_122}.
\end{minipage}
\label{tab:legal survive}
\end{table}

\newpage
%----------------------------------------------------------------------
\section{Tariff Implications for Indonesia}
%----------------------------------------------------------------------

Following the SCOTUS ruling, President Trump utilized Section~122 (\S122) and imposed global tariffs at 10\%, quickly raising them to 15\%---the statutory maximum under 19~USC~\S2132. For Indonesia, the critical question is what average tariff rate (ETR) its exporters now face relative to the negotiated deal.

Table~\ref{tab:etr} presents average tariff rates (ETR)---the trade-weighted average of applicable rates across all product categories---for five distinct policy scenarios tracing Indonesia's tariff trajectory. The figures use 2025 US import data from Census Bureau (\$34.84 billion total) and apply the Annex~II product exemptions that the Section~122 proclamation carries, which mirror the IEEPA Annex~II list.

\begin{table}[htbp]
\centering
\caption{Average Tariff Rates Across Scenarios}
\begin{tabular}{llcc}
\toprule
\textbf{Scenario} & \textbf{Period} & \textbf{ETR} & \textbf{Change from Baseline} \\
\midrule
1. Baseline & Pre-April 2025 & 5.0\% & -- \\
2. Liberation Day & April 2--9, 2025 & 36.6\% & +31.6 pp \\
3. Post-Deal & Aug 2025 -- Feb 19, 2026 & 19.7\% & +14.7 pp \\
4. Post-SCOTUS (10\% S122) & Feb 20, 2026 & 14.5\% & +9.5 pp \\
5. S122 at 15\% & Feb 21, 2026+ & 19.0\% & +14.0 pp \\
\bottomrule
\end{tabular}
\par\smallskip
\begin{minipage}{\linewidth}
\footnotesize
\textit{Note:} ETR (average tariff rate) is the trade-weighted average of applicable tariff rates across HTS 2-digit chapters, using 2025 U.S.\ import values as weights (\$34.84 billion). Applicable rates incorporate MFN duties, Section 232 tariffs, and the relevant reciprocal or Section 122 rate under each scenario. The stacking rule applies: the higher of Section 232 or reciprocal/S122 governs each product; tariffs are not additive. Scenarios 4 and 5 apply Annex II product exemptions (coffee, cocoa, partial electronics and rubber) as carried by the Section 122 proclamation. Change from baseline is in percentage points~(pp).\\
\textit{Source:} Authors' calculation based on Census Bureau 2025 trade data \cite{census2024}, USITC HTS \cite{usitc_hts}, WTO Tariff Profiles \cite{wto_tariff}, and Yale Budget Lab Tariff-ETR model \cite{yalebudgetlab_model}.
\end{minipage}
\label{tab:etr}
\end{table}

The headline finding is counterintuitive: at 15\%, Section~122 produces an average tariff rate for Indonesia (19.0\%) that is \textit{below}---not above---the rate Indonesia negotiated in the bilateral deal (19.7\%). The reason is structural. The Section~122 proclamation carries an Annex~II of exempt products that closely tracks the IEEPA Annex~II, covering coffee, cocoa, certain electronics components, and a subset of rubber products. Since those products were also exempted from the reciprocal tariff under the deal, the two regimes produce similar aggregate outcomes---with Section~122 coming in slightly lower because it also reduces tariffs on manufacturing sectors that faced higher deal-era rates.

The 5 percentage point escalation from 10\% to 15\% (between February 20 and 21) adds approximately \$1.55~billion in annual duties across all Indonesian exports. Every non-232, non-Annex~II sector is affected uniformly: apparel goes from 22\% to 27\%, footwear from 21\% to 26\%, and furniture from 11.5\% to 16.5\%. Only products covered by Section~232 (steel and aluminum at 50\%, automobiles at 25\%) are unaffected, because their 232 rates exceed the 15\% S122 rate and the stacking rule applies the higher of the two.

At 15\%, the administration has exhausted its unilateral authority under Section~122. The statute permits no further increase without congressional authorization. Combined with the 150-day time limit, this creates an unusual policy landscape: the tariff is simultaneously at its maximum level and approaching its expiration date. If Congress does not act before approximately July 20, 2026, non-232 tariffs would revert to the MFN baseline, dropping Indonesia's average tariff rate back toward 5.0\%.

\subsection{Winners and Losers: S122 at 15\% vs.\ the Deal}

Compared to the Post-Deal scenario, Section~122 at 15\% reshuffles the burden across Indonesian export sectors in ways that are not uniform. The story is less about the aggregate---which barely changes---and more about which sectors gain and which lose within that near-flat average.

Table~\ref{tab:sector impacts} presents sector-level average tariff rates across all five scenarios, with the final column showing each sector's position relative to the bilateral deal. The corrected figures apply Annex~II exemptions to the S122 scenarios and use 2025 import weights.

\begin{table}[htbp]
\centering
\caption{Sector-Level Average Tariff Rate Comparison}
\small
\begin{tabular}{lcccccccr}
\toprule
\textbf{Sector} & \textbf{HS} & \textbf{(\$B)} & \textbf{Base} & \textbf{Lib.} & \textbf{Deal} & \textbf{S122} & \textbf{S122} & \textbf{vs.} \\
& & & & \textbf{Day} & & \textbf{10\%} & \textbf{15\%} & \textbf{Deal} \\
\midrule
Electronics & 85 & 7.84 & 1.8\% & 33.8\% & 14.1\% & 10.7\% & 15.1\% & +1.0 pp \\
Footwear & 64 & 2.97 & 11.0\% & 43.0\% & 30.0\% & 21.0\% & 26.0\% & $-$4.0 pp \\
Knit Apparel & 61 & 2.52 & 12.0\% & 44.0\% & 29.1\% & 22.0\% & 27.0\% & $-$2.1 pp \\
Woven Apparel & 62 & 2.21 & 12.0\% & 44.0\% & 29.1\% & 22.0\% & 27.0\% & $-$2.1 pp \\
\textbf{Palm Oil} & \textbf{15} & \textbf{2.32} & \textbf{2.0\%} & \textbf{34.0\%} & \textbf{2.0\%} & \textbf{11.7\%} & \textbf{16.5\%} & \textbf{+14.5 pp} \\
\textbf{Rubber} & \textbf{40} & \textbf{1.71} & \textbf{3.0\%} & \textbf{35.0\%} & \textbf{3.0\%} & \textbf{10.3\%} & \textbf{14.0\%} & \textbf{+11.0 pp} \\
Furniture & 94 & 1.55 & 1.5\% & 33.5\% & 20.5\% & 11.5\% & 16.5\% & $-$4.0 pp \\
Machinery & 84 & 1.65 & 1.5\% & 33.5\% & 20.5\% & 11.5\% & 16.5\% & $-$4.0 pp \\
Coffee & 09 & 0.57 & 1.5\% & 33.5\% & 1.5\% & 1.5\% & 1.5\% & 0 pp \\
Cocoa & 18 & 0.77 & 3.0\% & 35.0\% & 3.0\% & 3.0\% & 3.0\% & 0 pp \\
Steel & 72--73 & 0.49 & 26.5\% & 34.0\% & 52.0\% & 52.0\% & 52.0\% & 0 pp \\
Aluminum & 76 & 0.21 & 28.0\% & 35.0\% & 53.0\% & 53.0\% & 53.0\% & 0 pp \\
\bottomrule
\end{tabular}
\par\smallskip
\begin{minipage}{\linewidth}
\footnotesize
\textit{Note:} Bold rows indicate sectors most affected by the shift from deal to S122. Import values (\$B) are 2025 full-year figures. ETR for each sector is the trade-weighted average tariff. Section~122 rates (Scenarios 4 and 5) apply Annex~II product exemptions: coffee and cocoa are fully exempt; rubber and electronics are partially exempt (27\% and 11\% of imports, respectively). ``vs.\ Deal'' is the difference between S122 at 15\% and Post-Deal ETRs. Automobile imports (HS~87, \$0.31B) face Section~232 at 25\%, which dominates the 15\% S122 rate.\\
\textit{Source:} Authors' calculation based on Census Bureau 2025 trade data, Yale Budget Lab Tariff-ETR model \cite{yalebudgetlab_model}, USITC HTS \cite{usitc_hts}, and ART deal exemption schedule \cite{art2026}.
\end{minipage}
\label{tab:sector impacts}
\end{table}

The most important---and perhaps surprising---finding in the sector table is the status of coffee and cocoa. Both sectors were exempted under Indonesia's bilateral deal. Both are also exempt under the Section~122 Annex~II. As a result, their tariff rates are essentially unchanged: coffee pays 1.5\% (MFN rate) under the deal and 1.5\% under Section~122; cocoa pays 3\% in both regimes. Indonesian coffee and cocoa exporters, accounting for approximately \$1.3 billion in 2025 exports, are unaffected by the regime change.

Palm oil and rubber tell a sharply different story. Neither commodity is covered by the IEEPA or S122 Annex~II; their zero-tariff treatment under the deal was Indonesia-specific---a concession wrested through months of negotiation, not a general product exemption. Under Section~122, both sectors lose that protection. Palm oil (\$2.32 billion, 85\% of US refined palm oil supply) faces a rate of 16.5\% instead of 2.0\% under the deal---an increase of 14.5 percentage points adding roughly \$336 million in annual duties. Rubber (\$1.71 billion, 48\% of US natural rubber supply) faces 14.0\% instead of 3.0\%---an increase of 11 percentage points adding approximately \$188 million. Combined, palm oil and rubber exporters pay roughly \$524 million more annually under Section~122 than they would have under the deal.

Manufacturing exporters, by contrast, benefit from the switch. Textiles and apparel (HS~61--62, \$4.73 billion, employing approximately 4 million workers directly) see their tariff drop from 29.1\% under the deal to 27\%---a 2.1 percentage point improvement. Footwear (\$2.97 billion) falls from 30\% to 26\%, saving roughly \$119 million annually. Furniture (\$1.55 billion) drops from 20.5\% to 16.5\%. These sectors were disadvantaged under the deal's 19\% blanket reciprocal rate; the flat Section~122 rate, being lower than 19\%, gives them modest relief.

Electronics (HS~85, \$7.84 billion---the single largest export category) occupies an intermediate position. The Section~122 Annex~II exempts approximately 11\% of HS~85 imports by value (primarily semiconductors and certain components), bringing the sector's average tariff to 15.1\% compared to 14.1\% under the deal. The marginal increase of 1 percentage point on a \$7.84 billion base translates to approximately \$78 million in additional duties---notable, but modest relative to the base.

Steel and aluminum remain in a category of their own. Section~232 tariffs at 50\% derive from separate legal authority, are unaffected by the SCOTUS ruling, and exceed the 15\% Section~122 rate, so the stacking rule leaves them unchanged at 52--53\%.

\subsection{The Legal Status of Indonesia's Concessions}

The Agreement of Reciprocal Trade signed on February 19, 2026 \cite{art2026} was negotiated under IEEPA authority that the Supreme Court invalidated the following day \cite{scotus2026}. This creates a fundamental asymmetry: Indonesia has made significant concessions whose value was predicated on tariff terms the US can no longer legally deliver.

Indonesia agreed to \$33 billion in US purchase commitments covering energy, Boeing aircraft, and agricultural products; elimination of tariffs on 99\% of US products exported to Indonesia; market access for US soybeans, corn, wheat, and cotton; and a 20-year extension of PT Freeport Indonesia's mining license from 2041 to 2061. These commitments were offered in exchange for a specific tariff structure---19\% reciprocal with product exemptions---that now has no statutory basis. International trade law principles of reciprocity suggest Indonesia has grounds to revisit or suspend its own commitments until the legal framework is clarified.

The Freeport MoU \cite{freeport_mou}, signed one day before the ruling, deserves particular scrutiny. Extending a major extractive concession in exchange for tariff terms that proved ephemeral raises questions about whether the agreement was structured as a standalone arrangement or as part of the broader trade package. Its enforceability may depend on that distinction.

\subsection{The Constraints of Section 122}

Section~122's design creates significant forward uncertainty for Indonesia. The 150-day clock started February 20, 2026, meaning the tariff expires around July 20, 2026, unless Congress authorizes an extension. This narrow window limits the tariff's coverage and creates a hard deadline for any legislative action.

More consequentially, Section~122 cannot differentiate by country. Indonesia now faces the same 15\% as Vietnam (which faced 46\% under IEEPA) or the European Union (which faced 20\%). The non-differentiation constraint eliminates the leverage that made bilateral deal-making possible. There is no mechanism under Section~122 for Indonesia to negotiate its way to a lower rate, because the statute does not permit country-specific rates.

At 15\%, the administration has exhausted its unilateral authority under this provision. Any further escalation requires an act of Congress---a higher bar in the current political environment.

%----------------------------------------------------------------------
\section{Forward-Looking Scenarios}
%----------------------------------------------------------------------

The internal dynamics of US trade policy will continue to have direct consequences for Indonesia's economy. The most likely near-term outcome is legislative action: Congress could pass new authority granting the president differentiated tariff power, potentially through a new statutory framework for reciprocal trade, an amendment to Section~301, or a standalone act. If Congress acts before the 150-day Section~122 expiry, the administration could reimpose country-specific rates. The level for Indonesia would depend heavily on the new legislation's structure and on whether Indonesia's concessions are treated as leverage or as already delivered.

The second scenario---Section~122 expiry without replacement---would benefit Indonesia substantially. Non-232 products would revert to MFN-only tariffs, dropping Indonesia's average tariff rate back toward 5.0\%. Steel and aluminum would remain at 50\% under Section~232, but the bulk of Indonesian exports would face only their pre-crisis baseline rates. This outcome would essentially reset the tariff situation to pre-Liberation Day levels, with Indonesia gaining the most among countries that faced high IEEPA rates.

A third possibility is deal restructuring under alternative legal authority. The administration could attempt to replicate the Indonesia deal using Section~301 (which requires a formal investigation into unfair trade practices), Section~201 (which requires an ITC injury finding, though typically on a sector-specific basis), or the rarely invoked Section~338 (authorizing retaliatory tariffs up to 50\% for discriminatory treatment). Any of these paths would take months and likely produce different tariff structures. Indonesia would be negotiating from a stronger position---not as a supplicant seeking relief from IEEPA rates, but as a country whose prior concessions were conditioned on terms now unavailable.

Finally, Treasury Secretary Bessent has indicated the administration may layer multiple authorities---232, 301, 122, 201, and 338---to maintain tariff pressure across product categories. This patchwork approach could maintain elevated tariffs on specific sectors while facing immediate legal challenges on the novel use of each authority. For Indonesia, the practical effect would be continued uncertainty about which products face which rates at any given time.

%----------------------------------------------------------------------
\section{Policy Recommendations}
%----------------------------------------------------------------------

The fast-changing tariff environment demands an equally agile response. For the Indonesian government, the first priority is strategic patience on the ART concessions. With the legal basis for the agreement invalidated, Indonesia should pause implementation of its own commitments---particularly tariff reductions on US products and agricultural market access---until the US clarifies what authority it will use and what tariff structure will actually be in place. Implementing concessions in exchange for tariff terms that no longer exist would be an asymmetric bargain.

The government should also develop explicit contingency plans for each of the four scenarios described above: congressional renewal, Section~122 expiry, restructuring under alternative authority, and authority layering. Engagement with US congressional leaders---not just the executive branch---is now essential, since Congress holds the key to whether the tariff landscape stabilizes or continues to shift.

The non-differentiation constraint of Section~122 is, paradoxically, a source of temporary advantage. Indonesia currently faces the same 15\% as Vietnam (which faced 46\% under IEEPA) and the EU (20\%). Indonesian exporters, particularly in apparel and footwear, have a narrow window to gain market share relative to higher-tariff competitors before Congress potentially restores country-differentiated rates. Export promotion efforts during the 150-day window should capitalize on this relative position.

On market diversification, the tariff instability reinforces a longer-term structural imperative. Deepening trade relationships with RCEP partners and advancing the CEPA negotiations with the EU reduces dependence on a single market governed by unpredictable policy. In the short run, diversification options are limited; in the medium term, reducing Indonesia's exposure to US tariff volatility is the most durable form of economic insurance.

For Indonesian businesses, the strategic advice divides by sector. Palm oil and rubber exporters face the most acute adjustment: a 14.5 and 11 percentage point tariff increase relative to the deal respectively, without any prospect of country-specific relief under Section~122. Building inventory, exploring US-customer cost-sharing arrangements, and monitoring Freeport-linked concession status are the immediate priorities. For the commodity sectors, Indonesia's dominant market share provides some pricing power---US buyers of refined palm oil (85\% sourced from Indonesia) and natural rubber (48\% share) have limited substitution options in the short run, which allows partial pass-through of the tariff cost to buyers.

Apparel, footwear, and furniture exporters, facing lower tariffs under Section~122 than under the deal, should accelerate shipments during the 150-day window. The 27\%, 26\%, and 16.5\% rates respectively may not last---congressional action could reimpose higher country-specific rates, or Section~122 could simply expire. The window of relative advantage is finite.

All exporters with goods that crossed US borders during the IEEPA period should consult customs brokers about duty drawback and refund possibilities. CBP protest deadlines are typically 180 days from liquidation; missing them forfeits potential refunds on tariffs that the SCOTUS ruling has now invalidated.

%----------------------------------------------------------------------
\section{Conclusions}
%----------------------------------------------------------------------

The ten months between Liberation Day and February 2026 compressed what would normally be years of trade policy evolution into a few headline-grabbing weeks. Indonesia signed a bilateral trade deal on February 19, watched it invalidated the next day, and absorbed a new tariff regime---all before the ink on the ART was dry. The practical tariff burden, correctly measured, turns out to be marginally better than what the deal would have provided in aggregate: 19.0\% under Section~122 at 15\% versus 19.7\% under the deal. But the aggregate masks a sharp redistribution of burden.

The corrected analysis---incorporating the Section~122 Annex~II product exemptions and 2025 trade data---changes both the numbers and the story. Coffee and cocoa exporters, initially reported as major losers from the regime change, are in fact Annex~II exempt under both IEEPA and Section~122, leaving their tariff treatment essentially unchanged from the deal. Palm oil and rubber exporters bear the genuine losses: their deal-specific exemptions have no equivalent in the S122 Annex~II, and the tariff increases of 11--15 percentage points relative to the deal translate to roughly \$524 million in additional annual duties. Manufacturing sectors---apparel, footwear, furniture---gain modest relief.

The deeper problem is structural uncertainty. Section~122 expires around July 2026, the US has no current legal authority to impose country-differentiated tariffs, and the legislative path to new authority is uncertain. Indonesia's concessions under the ART were made in exchange for tariff terms that can no longer be delivered. Pausing those concessions, engaging Congress directly, and using the current window of relative advantage in manufacturing exports are the immediate priorities. The longer-term lesson is clear: any trade relationship with the United States must now be built on authorities that the Supreme Court has validated, not on executive orders whose legal standing is contested.

%----------------------------------------------------------------------
\section*{Appendix 1: Legal Contexts}
%----------------------------------------------------------------------

\subsection*{The SCOTUS Decision: \textit{Learning Resources, Inc.\ v.\ Trump}}

In a 6--3 decision authored by Chief Justice Roberts, the Supreme Court held that ``IEEPA does not authorize the president to impose tariffs.'' The majority---Roberts, Sotomayor, Kagan, Gorsuch, Barrett, and Jackson---reasoned that IEEPA contains no reference to tariffs or duties; Congress has never used the word ``regulate'' to authorize taxation; no president had previously interpreted IEEPA as conferring tariff power; and the Constitution grants tariff authority to ``Congress alone'' (Art.\ I, Sec.\ 8). Thomas and Kavanaugh (joined by Alito) dissented.

\subsection*{The Section 122 Pivot and Escalation}

Within hours of the ruling, the Trump administration invoked Section~122 of the Trade Act of 1974 \cite{s122_statute,crs_122}, initially imposing a 10\% global tariff, raised to 15\% the next day. Key constraints: the current rate of 15\% is the statutory ceiling; the maximum duration is 150 days without congressional authorization; rates must apply uniformly without country differentiation; the statutory trigger is a presidential determination of ``large and serious'' balance-of-payments deficit; and at 15\%, the administration has exhausted its unilateral authority under this statute. This was the first use of Section~122 to impose tariffs in the statute's history.

%----------------------------------------------------------------------
\section*{Appendix 2: Calculating the Average Tariff Rate for Indonesia}
%----------------------------------------------------------------------

\subsection*{Methodology}

The average tariff rate (ETR, effective tariff rate) is calculated as:

\begin{equation}
\text{ETR} = \frac{\sum_{i} (\text{applicable\_rate}_i \times \text{import\_value}_i)}{\sum_{i} \text{import\_value}_i}
\end{equation}

\noindent where $i$ denotes each HTS-2 chapter. The applicable rate for each chapter incorporates the MFN tariff rate (baseline duty from the Harmonized Tariff Schedule \cite{usitc_hts,wto_tariff}), Section~232 tariffs (steel and aluminum at 50\%, automobiles at 25\%), and the relevant reciprocal or Section~122 rate under each scenario.

The stacking rule (per Yale Budget Lab methodology \cite{yalebudgetlab_model}): Section~232 and reciprocal/Section~122 tariffs are mutually exclusive; the higher additional rate applies. Products are not double-charged.

For Scenarios 4 and 5 (Section~122), Annex~II product exemptions are applied. The exempt fraction for each HTS-2 chapter is estimated from the Yale Budget Lab's 12-4 scenario output: chapters with below-average ETR in that scenario have a proportional share of Annex~II-exempt imports. Specifically, exempt fraction $\alpha_c = \max(0, 1 - \text{ETR}_{c}^{12\text{-}4}/0.19)$ for chapters where $\text{ETR}_{c}^{12\text{-}4} \leq 0.19$; 232-dominated chapters ($\text{ETR}_{c}^{12\text{-}4} > 0.19$) receive no correction. Import weights use 2025 US Census Bureau data (\$34.84 billion). The 2025 data reflects significant pre-tariff front-loading, particularly in electronics.

\subsection*{Scenario Definitions}

Five scenarios trace the evolution of US tariff policy toward Indonesia from the pre-crisis baseline through the SCOTUS ruling and the Section~122 escalation.

\textbf{Scenario 1: Baseline (Pre-April 2025).}
MFN rates only for most products; Section~232 at 25\% on steel (HS~72--73) and aluminum (HS~76); no reciprocal tariffs.

\textbf{Scenario 2: Liberation Day (April 2--9, 2025).}
32\% IEEPA reciprocal tariff on all Indonesian products; Section~232 at 25\% superseded by higher 32\% reciprocal; effectively MFN~+ 32\% across the board.

\textbf{Scenario 3: Post-Deal (August 2025 -- February 19, 2026).}
19\% reciprocal tariff (reduced from 32\%); 1,819 tariff lines exempted at 0\% reciprocal (palm oil, rubber, coffee, cocoa, spices, semiconductor components); textile TRQ ($\sim$10\% of HS~61--62 trade at MFN-only); Section~232 at 50\% on steel/aluminum; Section~232 at 25\% on autos.

\textbf{Scenario 4: Post-SCOTUS at 10\% S122 (February 20, 2026).}
10\% Section~122 tariff, flat and global; Annex~II product exemptions applied (coffee, cocoa, partial electronics and rubber); Section~232 at 50\% on steel/aluminum (higher than 10\%, so 232 applies); Section~232 at 25\% on autos.

\textbf{Scenario 5: S122 at 15\%---Statutory Maximum (February 21, 2026+).}
15\% Section~122 tariff---the statutory ceiling under 19~USC~\S2132; Annex~II product exemptions applied; Section~232 at 50\% on steel/aluminum; Section~232 at 25\% on autos.

%----------------------------------------------------------------------
\section*{Appendix 3: Data Sources and Validation}
%----------------------------------------------------------------------

\subsection*{Cross-Validation}

Import data are from the US Census Bureau monthly trade statistics for 2025, aggregated by HTS 2-digit chapter for Indonesia (country code 5600). The 2025 full-year total of \$34.84 billion compares to \$28.05 billion in 2024---a 24.2\% increase consistent with widespread reporting of pre-tariff import front-loading. Tariff rates are cross-validated against the Yale Budget Lab's open-source Tariff-ETR model \cite{yalebudgetlab_model}:

The Yale baseline US ETR is 2.4\% (January 2025). Indonesia's baseline of 5.0\% is higher due to its concentration in textiles and footwear, which carry above-average MFN rates---consistent with expectations. The Yale post-SCOTUS US ETR is 9.1\% at 10\% S122 without Annex~II correction. Indonesia's 14.5\% at 10\% S122 (with Annex~II) reflects the higher MFN rates on its product mix and the partial application of exemptions. The Yale Liberation Day increase is +11.5~pp overall; Indonesia's jump to 36.6\% is consistent with a 32\% tariff pushing a country with a relatively high baseline.

\subsection*{Reconciliation with Yale Budget Lab Model (Incremental ETR)}

This report's ETR includes MFN baseline duties (total tariff burden). The Yale Budget Lab model measures incremental ETR---the additional burden from executive actions above the MFN schedule. The difference reflects the trade-weighted average MFN rate for Indonesian goods ($\sim$2.6~pp):

\begin{table}[htbp]
\centering
\caption{ETR Reconciliation: Total vs.\ Incremental}
\begin{tabular}{lcccl}
\toprule
\textbf{Scenario} & \textbf{Total ETR} & \textbf{Incremental ETR} & \textbf{Diff.} & \textbf{Explanation} \\
\midrule
Post-Deal / 12-4 & 19.7\% & 18.0\% & $\sim$1.7 pp & MFN adds $\sim$5 pp; deal \\
& & & & exemptions subtract $\sim$3 pp \\
S122 at 10\% / 2-20 & 14.5\% & 11.9\% & $\sim$2.6 pp & Trade-weighted avg MFN \\
S122 at 15\% / 2-21 & 19.0\% & 16.4\% & $\sim$2.6 pp & Same MFN gap; Annex II \\
& & & & applied in both calcs \\
\bottomrule
\end{tabular}
\par\smallskip
\begin{minipage}{\linewidth}
\footnotesize
\textit{Note:} Total ETR includes MFN baseline duties. Incremental ETR (Yale model) measures only executive-action tariffs above MFN. The consistent $\sim$2.6~pp gap in the S122 scenarios reflects the trade-weighted average MFN rate for Indonesian products.\\
\textit{Source:} Total ETR: authors' calculation; Incremental ETR: authors' calculation using Yale Budget Lab model \cite{yalebudgetlab_model}.
\end{minipage}
\end{table}

Total ETR answers: ``What is the total duty rate on Indonesian goods entering the US?'' Incremental ETR answers: ``How much additional burden do executive actions impose beyond the MFN baseline?'' Both measures are valid and complementary. For trade cost analysis and exporter decision-making, the total ETR reflects actual duty payments.


\begin{thebibliography}{99}

\bibitem{yalebudgetlab2026}
Yale Budget Lab, ``State of U.S.\ Tariffs,'' February 20, 2026.
\url{https://budgetlab.yale.edu/research/state-us-tariffs}

\bibitem{yalebudgetlab_model}
Yale Budget Lab, \textit{Tariff-ETRs}: Open-Source Effective Tariff Rate Model, GitHub.
\url{https://github.com/Budget-Lab-Yale/Tariff-ETRs}

\bibitem{yalebudgetlab2025jul}
Yale Budget Lab, ``State of U.S.\ Tariffs,'' July 23, 2025.

\bibitem{yalebudgetlab2025nov}
Yale Budget Lab, ``State of U.S.\ Tariffs,'' November 17, 2025.

\bibitem{census2024}
U.S.\ Census Bureau, \textit{Foreign Trade --- U.S.\ Trade in Goods with Indonesia}, 2024--2025.
\url{https://www.census.gov/foreign-trade/balance/c5600.html}

\bibitem{ustr2024}
Office of the United States Trade Representative (USTR), Indonesia Country Commercial Guide / Trade Profile, 2024.

\bibitem{usitc_hts}
U.S.\ International Trade Commission (USITC), \textit{Harmonized Tariff Schedule of the United States}, 2025 revision.

\bibitem{wto_tariff}
World Trade Organization, \textit{World Tariff Profiles --- United States}, 2024/2025.

\bibitem{crs_232}
Congressional Research Service, ``Section 232 Investigations: Overview and Issues for Congress,'' IN12519.

\bibitem{crs_122}
Congressional Research Service, ``Section 122 of the Trade Act of 1974,'' R48435.

\bibitem{scotus2026}
\textit{Learning Resources, Inc.\ v.\ Trump}, 607 U.S.\ \_\_\_ (2026).

\bibitem{s122_statute}
Trade Act of 1974, \S 122, codified at 19 U.S.C.\ \S 2132.

\bibitem{art2026}
Agreement of Reciprocal Trade (ART), United States--Indonesia, signed February 19, 2026.

\bibitem{freeport_mou}
Memorandum of Understanding on the Extension of PT Freeport Indonesia's Mining License, signed February 19, 2026.

\bibitem{comtrade}
United Nations COMTRADE Database, via Trading Economics, U.S.--Indonesia bilateral trade data.

\bibitem{usitc_dataweb}
USITC DataWeb, HTS-level U.S.\ import statistics.
\url{https://dataweb.usitc.gov}

\end{thebibliography}

\end{document}
